% =====================================================================
%  PROTOCOLO DE SESIÓN DE LA SECUENCIA (hoja de campo del tutor)
%  Una copia por sesión. Sirve para cualquiera de las cuatro tareas.
%  Fundamentación de las tareas: secuencia_tareas.tex
%  Compilar con: pdflatex protocolo_sesion.tex (dos veces)
% =====================================================================
\documentclass[10pt,letterpaper]{article}

\usepackage[utf8]{inputenc}
\usepackage[T1]{fontenc}
\usepackage[spanish,es-tabla,es-nodecimaldot]{babel}
\usepackage{lmodern}
\usepackage{microtype}
\usepackage[margin=1.5cm,top=1.8cm]{geometry}
\usepackage{amsmath,amssymb}
\usepackage{tabularx,array}
\usepackage{enumitem}
\usepackage[table]{xcolor}
\usepackage[most]{tcolorbox}
\usepackage{fancyhdr}
\usepackage{lastpage}
\usepackage{url}

\definecolor{unalverde}{RGB}{0,98,65}
\definecolor{grisclaro}{RGB}{242,242,242}

\newcolumntype{L}{>{\raggedright\arraybackslash}X}
\setlist{itemsep=1pt,topsep=3pt}
\setlength{\parindent}{0pt}

\newtcolorbox{guion}{colback=white,colframe=unalverde!60,leftrule=4pt,boxrule=0.6pt,arc=0pt,fontupper=\itshape,left=6pt,top=2pt,bottom=2pt,before skip=3pt,after skip=3pt}
\newcommand{\momento}[2]{\par\smallskip{\color{unalverde}\large\bfseries #1}\hfill{\small #2}\par\vspace{2pt}}
\newcommand{\cas}{$\square$\;}

\pagestyle{fancy}
\fancyhf{}
\fancyhead[L]{\small Protocolo de sesión --- Secuencia de tareas con Gemini}
\fancyhead[R]{\small Código: \rule{1.4cm}{0.4pt} \quad Tarea: T\rule{0.5cm}{0.4pt} \quad Fecha: \rule{1.8cm}{0.4pt}}
\fancyfoot[C]{\small \thepage\ de \pageref{LastPage}}
\renewcommand{\headrulewidth}{0.4pt}

\begin{document}

{\Large\bfseries Protocolo de sesión\par}
{\small Una tarea por sesión · 55--60 min de trabajo · Tiempos provisionales, se ajustan en el pilotaje\par}
\medskip

\begin{tcolorbox}[colback=grisclaro,colframe=unalverde,arc=2pt,title=\textbf{Antes de que llegue el estudiante (5 min)},fonttitle=\small]
\small
\cas Consentimiento firmado (y autorización de audio: sí / no) \quad
\cas Hoja de la tarea impresa \quad
\cas Esferos azul o negro y rojo \quad
\cas Grabadora lista \\
\cas Gemini abierto en la cuenta institucional del investigador, \textbf{conversación nueva} \quad
\cas Modo \textbf{aprendizaje guiado} activado \quad
\cas Instrucción inicial pegada (abajo) \quad
\cas GeoGebra disponible
\end{tcolorbox}

\textbf{Instrucción inicial para Gemini} (igual para todos; no lleva datos del estudiante):
\begin{guion}
Soy estudiante de un primer curso de cálculo en la universidad. Estoy resolviendo una tarea sobre límites y ya intenté resolverla por mi cuenta. Quiero entender, no que me des la respuesta: guíame con preguntas y pistas, y si cometo un error, ayúdame a encontrarlo en lugar de corregirlo directamente.
\end{guion}

\textbf{Problema de la tarea} (taller de origen y número): \rule{9cm}{0.4pt}

\momento{Apertura}{2 min}
\begin{guion}
Hoy vamos a trabajar un ejercicio en tres partes. Primero solo, con el esfero azul. Después con Gemini, y lo que cambie lo escribe en rojo sin borrar lo azul. Al final decide cuál es su respuesta y por qué. Me interesa cómo piensa, no solo el resultado. Trate de decir en voz alta lo que va pensando.
\end{guion}

\momento{Momento 1 --- Sin IA, esfero azul o negro}{20 min}
Entregar la hoja. El tutor \textbf{no explica}. Si hay bloqueo de dos minutos o el estudiante pide ayuda, usar la escala en orden y anotar el nivel en el registro:

\begin{center}\small
\begin{tabularx}{\linewidth}{@{}lL@{}}
\textbf{A0} Silencio & Esperar 10 segundos.\\
\textbf{A1} Explicitación & ¿Qué ha intentado hasta ahora?\\
\textbf{A2} Reformulación & ¿Cómo diría con sus palabras lo que se pide?\\
\textbf{A3} Heurística & ¿Le serviría una gráfica, una tabla o un dibujo?\\
\textbf{A4} Contenido & Se da el paso y se pide que lo explique.\\
\end{tabularx}
\end{center}
Preguntas para sostener el razonamiento: ¿Por qué hizo ese paso? · ¿Cómo sabe que eso es cierto? · ¿Hay otra manera de verlo?

\momento{Momento 2 --- Con Gemini, esfero rojo, sin borrar}{20 min}
\begin{guion}
Ahora puede usar Gemini. Antes de enviar cada mensaje, dígame qué le va a preguntar y para qué. Si cambia algo en la hoja, hágalo en rojo, sin borrar.
\end{guion}
\begin{itemize}
  \item Propósito previsto de Gemini en esta tarea: \cas contrastar (T1) \quad \cas explicar (T2) \quad \cas verificar (T3, T4)
  \item Si Gemini entrega la solución completa, no se interrumpe: se anota y se pregunta después ``¿qué hizo con esa respuesta?''.
  \item Con Gemini abierto, el tutor no da ayudas A3 ni A4.
\end{itemize}

\momento{Momento 3 --- Validación}{10 min}
\begin{itemize}
  \item ¿Cuál es su respuesta final? Escríbala y explique por qué está seguro.
  \item Compare lo azul con lo rojo: ¿qué cambió y qué lo hizo cambiar?
  \item ¿Qué parte de la resolución hizo usted y qué parte hizo Gemini?
\end{itemize}

\momento{Cierre}{2 min}
\begin{guion}
Gracias. Lo que escribió no afecta su nota. La próxima vez seguimos con la siguiente tarea.
\end{guion}

\begin{tcolorbox}[colback=white,colframe=orange!70!black,arc=2pt,title=\textbf{Para que la tarea no pierda su exigencia},fonttitle=\small]
\small
Henningsen y Stein (1997): la demanda de una tarea se mantiene cuando hay presión sostenida por explicar, tiempo suficiente y andamiaje que no le quita la complejidad, y cae cuando alguien ``toma el control'' de las partes difíciles o cuando la atención pasa del significado a la respuesta correcta.\\[2pt]
\textbf{En la sesión:} no resolver por el estudiante · no cortar el momento 1 antes de tiempo · preguntar ``por qué'' aunque la respuesta esté bien · tratar el error como punto de partida.
\end{tcolorbox}

% ---------------------------------------------------------------------
\newpage
\section*{Registro de la sesión \small(completar en los 30 minutos siguientes)}
\renewcommand{\arraystretch}{1.5}
\begin{tabularx}{\linewidth}{|p{5.4cm}|X|}
\hline
Duración real de cada momento (M1 / M2 / M3) & \\ \hline
Estrategia inicial (M1) & \\[0.8cm] \hline
Representaciones usadas sin que se pidieran & \\ \hline
Bloqueos (minuto y qué pasó) & \\[0.8cm] \hline
Ayudas dadas (nivel y minuto) & \\ \hline
Nivel máximo de ayuda & \cas A0 \quad \cas A1 \quad \cas A2 \quad \cas A3 \quad \cas A4 \\ \hline
Número de mensajes a Gemini & \\ \hline
Qué le preguntó a Gemini y para qué & \\[0.8cm] \hline
Qué hizo con las respuestas & \cas Las integró \quad \cas Las contrastó \quad \cas Pidió aclaración \quad \cas Las ignoró \quad \cas Copió \\ \hline
¿Gemini entregó la solución completa? & \cas No \quad \cas Sí (en qué mensaje): \\ \hline
Cambios de azul a rojo & \\[0.8cm] \hline
Respuesta final y justificación (M3) & \\[0.8cm] \hline
Qué atribuye a sí mismo y qué a Gemini & \\[0.8cm] \hline
Errores del análisis a priori que aparecieron & \\ \hline
Errores no previstos & \\ \hline
¿Induje alguna respuesta? ¿Dónde? & \\ \hline
Incidentes técnicos & \\ \hline
Ajuste para la tarea o el protocolo & \\[0.8cm] \hline
\end{tabularx}

\bigskip
\textbf{Después de la sesión}
\begin{itemize}
  \item[\cas] Copiar la conversación de Gemini a \texttt{E\#\#\_T\#\_gemini.txt} (sin usar ``Compartir''), reemplazar cualquier dato identificable por el código y \textbf{borrarla de la cuenta}.
  \item[\cas] Escanear la hoja \textbf{a color} como \texttt{E\#\#\_T\#\_hoja.pdf}.
  \item[\cas] Guardar audio y archivos en \texttt{datos\_investigacion/} (fuera del repositorio).
  \item[\cas] Transcribir con el comando \texttt{codificador} en menos de 72 horas.
\end{itemize}

\medskip
{\footnotesize Henningsen, M., y Stein, M. K. (1997). Mathematical tasks and student cognition: Classroom-based factors that support and inhibit high-level mathematical thinking and reasoning. \emph{Journal for Research in Mathematics Education, 28}(5), 524--549. \url{https://doi.org/10.5951/jresematheduc.28.5.0524}}

\end{document}
